%% kapa: exam=1 round=37 year=2026 period=1 form=A subject=economics subject_ko=경제학원론 seq=13 number=53
%% source: Q-Net 감정평가사 기출문제 (cst00309) artlSeq=5258924 file=raw/1008/5258924/2026년 제37회 감정평가사 1차 1교시.pdf
%% conversion: kordoc 4.24.51 -> tools/gen1.py (auto); review=proofread-sonnet; flags=none
\begin{question}{53}
소득-여가 결정모형에서 효용극대화를 추구하는 甲의 노동공급에 관한 설명으로 옳은 것을 모두 고른 것은? (단, 소득과 여가는 모두 정상재이며, 소득효과 및 대체효과의 크기 비교는 절댓값을 기준으로 한다.)
\begin{framedtext}
ㄱ. 시간당 임금이 상승할 경우, 대체효과는 노동공급 감소요인이다. \par
ㄴ. 시간당 임금이 하락할 경우, 소득효과가 대체효과보다 크다면 노동공급은 증가한다. \par
ㄷ. 시간당 임금이 상승할 경우, 소득효과 및 대체효과의 크기와 상관없이 노동공급은 증가한다. \par
\end{framedtext}
\begin{choices}
  \item ㄱ
  \item ㄴ
  \item ㄷ
  \item ㄱ, ㄴ
  \item ㄴ, ㄷ
\end{choices}
\end{question}
\answer{2}
